This section assesses specific vulnerabilities in PostgreSQL/pgvector-based memory architectures, with particular attention to systems using automated conversation capture (``autoprocessors'').

\subsection{Architecture Under Analysis}

The reference architecture consists of:

\begin{itemize}
    \item \textbf{Storage}: PostgreSQL with pgvector extension for embedding storage
    \item \textbf{Embedding Model}: Sentence transformer (e.g., e5-base-v2) for text-to-vector conversion
    \item \textbf{Autoprocessor}: Background service that captures conversation exchanges and stores them as memories
    \item \textbf{Retrieval}: Cosine similarity search over embeddings, typically returning top-$k$ results
    \item \textbf{Metadata}: Tags, significance scores, timestamps associated with each memory
\end{itemize}

This architecture is representative of production AI agent memory systems and is, in fact, the architecture underlying the author's own cognitive infrastructure.

\subsection{Vulnerability 1: Autoprocessor as Injection Vector}

The autoprocessor creates a direct pathway from conversation content to persistent memory. This is the primary attack surface.

\begin{securitywarning}
Any content that appears in a conversation---including malicious instructions, poisoned demonstrations, or attack payloads---can be captured and persisted as a ``legitimate'' memory. The autoprocessor does not distinguish between user input, agent output, or potentially adversarial content.
\end{securitywarning}

\textbf{Attack Scenario}: An adversary engages the agent in conversation, gradually introducing poisoned reasoning patterns. The autoprocessor captures these exchanges. On subsequent sessions, when the agent retrieves contextually relevant memories, it may retrieve the poisoned conversation and imitate the adversarial patterns.

\textbf{Risk Level}: \textcolor{red}{\textbf{HIGH}}

\subsection{Vulnerability 2: Provenance Blindness}

Standard pgvector schemas do not track who or what generated each embedding. A memory retrieved from the store has no cryptographic attestation of its origin.

\textbf{Consequences}:
\begin{itemize}
    \item Cannot distinguish self-generated memories from externally-injected content
    \item Cannot verify that a memory was created by a trusted process
    \item Cannot detect tampering or modification of existing memories
\end{itemize}

\textbf{Risk Level}: \textcolor{red}{\textbf{HIGH}}

\subsection{Vulnerability 3: Similarity Without Safety}

Retrieval is based purely on semantic similarity. A memory that is highly relevant to the current task will be retrieved regardless of whether its content is safe to imitate.

\textbf{Attack Implication}: Poisoned entries can be designed to maximize relevance to high-value task categories (e.g., data processing, code execution), ensuring they are retrieved precisely when they can cause the most harm.

\textbf{Risk Level}: \textcolor{orange}{\textbf{MEDIUM-HIGH}}

\subsection{Vulnerability 4: Shared Memory Pools}

If multiple agents or users share a memory store, poisoning by one party affects all others. This is particularly concerning in multi-agent systems where agents may have different trust levels.

\textbf{Risk Level}: \textcolor{red}{\textbf{HIGH}} (if applicable) / \textcolor{green}{\textbf{LOW}} (if single-agent)

\subsection{Vulnerability 5: Retrieval Parameter Sensitivity}

Research shows that attack success rates correlate with retrieval aggressiveness. With $k=3$, GPT-4o-mini showed 6\% attack success rate; with $k=10$, this rose to 38\% \cite{memorypoisoning2025}.

\textbf{Implication}: Systems configured for more comprehensive memory retrieval are more vulnerable.

\textbf{Risk Level}: \textcolor{orange}{\textbf{MEDIUM}}

\subsection{Vulnerability 6: Lack of Temporal Decay}

Without temporal weighting, old memories have equal retrieval probability to recent ones. This creates a long attack window---poison injected months ago remains fully active.

\textbf{Risk Level}: \textcolor{yellow}{\textbf{LOW-MEDIUM}}

\subsection{Summary: Attack Surface Assessment}

\begin{table}[h]
\centering
\caption{Vulnerability Assessment Summary}
\begin{tabular}{lcp{6cm}}
\toprule
\textbf{Vulnerability} & \textbf{Risk} & \textbf{Mitigation Priority} \\
\midrule
Autoprocessor injection & HIGH & Immediate: pre-storage screening \\
Provenance blindness & HIGH & High: cryptographic attestation \\
Similarity without safety & MEDIUM-HIGH & High: safety reranking \\
Shared memory pools & HIGH/LOW & Context-dependent: RLS \\
Retrieval sensitivity & MEDIUM & Medium: parameter tuning \\
No temporal decay & LOW-MEDIUM & Medium: trust decay \\
\bottomrule
\end{tabular}
\label{tab:vulnerability-summary}
\end{table}
